\section{Constructed causal labels, semantic refusal and the unit of evaluation}
\label{exec17:benchmark-development}

The preceding chapter established what an evaluation must preserve. This chapter asks which simpler or more established method supplies a credible comparison at each stage. Constructed labels and literature retrieval test the reference; coding and repair studies test calculation; factor extraction and simulated respondents test whether the inputs represent the intended population. The final program-of-thought study returns to the full chain, where saved accuracy, scientific meaning and executable reuse can be assessed separately.

A causal assistant can improve the wording of an answer while leaving its
decision content unchanged. It can also repeat a familiar causal story more
fluently than an analyst without establishing that the story applies to the
bank's observations. Evaluation therefore needs to distinguish recognizing a
stipulated mechanism, identifying an effect from supplied information and
choosing a useful action. Two further benchmarks make this distinction
concrete: the published CELLO study by Chen and colleagues and the first
Research Square version of CausalReasonBench by Cherif. The former uses visual
scenes and constructed graphs; the latter uses textual scenarios and automatic
reasoning scores. Their results concern their respective tasks and historical
model settings. Neither is a measurement of a causal banking
workflow.~\citep{exec17_benchmark_cello,exec17_benchmark_cherif}

Consider the proposed deposit-pricing assistant. A manager asks whether a
retention offer will prevent withdrawals. A useful response must distinguish
the association between receiving an offer and staying with the bank from the
effect of offering the same terms to an eligible customer. It must also state
what the information cannot establish. A benchmark that supplies a complete
causal mechanism tests whether the assistant follows that mechanism. A
benchmark that supplies only an observational vignette may instead need an
answer that preserves several possible mechanisms. Those are different
questions even when the surface wording and requested answer are identical.
Constructed labels can be valuable for the first question; their construction
does not resolve the second.

\subsection{What a scene-derived causal answer establishes}

CELLO builds causal reasoning questions from Visual Genome relations. The
construction matches subgraphs to graph templates and uses language-model
generation to produce questions and distractors. Its prompting procedure
separates entity extraction, graph construction, task recognition and causal
knowledge. This organization is relevant to the proposal because a banking
assistant also needs to turn entities and relationships into an explicit
question before asking a numerical engine to answer it. The distinction is
that a relation visible in an image, or recorded in a customer document, need
not determine the intervention semantics assigned by a template. Static
co-occurrence, a plausible physical explanation and a warranted counterfactual
are distinct kinds of information. The benchmark's annotation and image
limitations therefore constrain the interpretation of its scores, rather than
being mere implementation details.~\citep{exec17_benchmark_cello}

The paper reports a test set of 1,200 questions, with 100 selected for each of
twelve tasks. Its Appendix B.2 describes a separate validity check of 120
instances, ten per task, by two annotators: 91.7\% were judged valid, with
reported Cohen agreement of 84.1\%. These figures describe a question-validity
review. They do not turn each template's intervention or counterfactual claim
into experimentally established truth. For banking, the analogous review
would check that the prompt faithfully expresses an agreed mechanism and
information set. It would still leave open whether that mechanism describes
the actual offer assignment, customer selection and outcome clock.
\citep{exec17_benchmark_cello}

Table 2 reports GPT-4o overall accuracy increasing from 0.59 to 0.70 with
CELLO-CoT, while confounder identification increases from 0.29 to 0.35.
Task-specific behavior matters: an aggregate improvement does not warrant
using a system to select adjustment variables. There is also an unresolved
arithmetic boundary in the published table. The twelve displayed
GPT-4o-plus-CELLO-CoT task entries are
\[
 0.83,\ 0.70,\ 0.52,\ 0.50,\ 0.35,\ 0.75,
 \quad 0.81,\ 0.72,\ 0.79,\ 0.77,\ 0.90,\ 0.61.
\]
Their equal-task mean is 0.6875. Under conventional rounding to two decimal
places, even a conservative closed envelope for that mean is
\begin{equation}
 [0.6825,\,0.6925]\ \cap\ [0.695,\,0.705]=\varnothing,
 \label{exec17:benchmark-rounding}
\end{equation}
where the second interval is the corresponding envelope for the reported
overall 0.70. With exactly 100 binary correctness outcomes per task, the
displayed entries would sum to 825 correct answers out of 1,200. This is a
conditional aggregation check against the published original. The cause of
the difference is unknown; no sample-level result or alternative aggregation
has been reconstructed. Consequently, the reported improvement remains an
author report with this unresolved qualification, rather than an independently
reproduced effect of the prompting procedure.
\citep{exec17_benchmark_cello,exec17_benchmark_exact}

\subsection{Refusal can be correct for the wrong reason}

CELLO's robustness exercise reformulates the questions as requests with
additional premises and binary answers. Section 6.4 states that the correct
answer is consistently ``No.'' This is a separate distribution from the
nearly balanced original binary questions. A constant-No classifier scores
100\% on the reformulated test without inspecting the scene. The paper also
describes GPT-4o declining requests because it is a language model unable to
interact with the physical world, rather than explaining the scene's causal
constraints. Correct answer polarity can therefore conceal a mismatch
between the tested mechanism and the reason for refusal.
\citep{exec17_benchmark_cello}

For a banking assistant, ``I cannot execute a payment'' and ``the available
evidence does not identify the effect of this payment policy'' have different
meanings. The first states an action limitation; the second states an
information limitation. A system could utter either sentence regardless of
whether a proposed comparison is scientifically valid. An evaluation should
therefore record the answer, the asserted reason and the consequence for the
requested decision. It should include supported requests as well as
unsupported requests, and distinguish the two kinds of refusal. Otherwise a
system that rejects every causal analysis may appear dependable while
providing no usable analysis.

Publicly constructed paired cases offer an inexpensive way to test this
distinction before commissioning a bank experiment. Begin with a mechanism
whose query is identified under the supplied assumptions. Create a companion
case that removes the assignment information or introduces an unresolved
alternative mechanism while keeping the business language similar. The
supported case should receive the warranted query or bounded answer; the
unsupported case should preserve uncertainty and name the missing information.
A complementary pair can retain the information and change only whether the
assistant has authority to act. This tests whether the assistant confuses
execution authority with causal identification. These are proposed evaluation
contrasts. Their merit is that the mechanism and the changed information are
explicit; a public observational label alone would not supply their causal
truth.

\subsection{Lexical agreement and directional reasoning}

CausalReasonBench supplies twelve underlying scenarios spanning four domains
and three difficulties. It repeats them over four task types, four prompting
strategies, three models and two replicates, yielding 1,152 responses, or 384
per model. The models are the paper's Groq-hosted LLaMA 3.1-8B, LLaMA 3.3-70B
and Qwen3-32B settings. The main settings use temperature zero, with the initial
contrast-case pass at 0.1 and seed 42 where supported. These details delimit a
historical evaluation; the response count is not a count of independently
sampled causal mechanisms. The gold explanations are author-defined using
the stated reference sources. In the minimum-wage example, the appendix itself
notes empirical debate. Following a stipulated textbook mechanism is a
different task from resolving that debate from an observational
vignette.~\citep{exec17_benchmark_cherif}

The paper's Causal Identification Rate, abbreviated CIR, is a token-set F1
score. For response-token set $R$, gold-keyword set $G$ and
$\varepsilon=10^{-8}$, its precision and recall are
$P=|R\cap G|/(|R|+\varepsilon)$ and
$Q=|R\cap G|/(|G|+\varepsilon)$, and
$\mathrm{CIR}=2PQ/(P+Q+\varepsilon)$. Causal Logic Precision, CLP, separately
assigns 1 to a gold forward pattern, 0.5 to ambiguous directional language,
and 0 to a reverse or absent pattern. Counterfactual Coherence Ratio, CCR,
uses counterfactual-token F1 for counterfactual tasks and equals CIR otherwise.
Thus the composite score outside the counterfactual tasks reduces to
\begin{equation}
 \mathrm{CAR}
 =0.40\mathrm{CIR}+0.35\mathrm{CLP}+0.25\mathrm{CCR}
 =0.65\mathrm{CIR}+0.35\mathrm{CLP}.
 \label{exec17:benchmark-score}
\end{equation}
The apparent three-component assessment has two distinct components in that
regime. This follows directly from the published definitions; it is not an
empirical claim about the models.
\citep{exec17_benchmark_cherif,exec17_benchmark_exact}

Set overlap ignores token order. For example, ``funding causes deposits'' and
``deposits causes funding'' give the same token set and hence the same CIR
against that gold set. The deliberately terse reverse statement isolates the
set operation; grammatical correction is a separate language task. CLP is
intended to test direction and could distinguish the two. The paper describes
whitespace tokenization, case normalization, Porter stemming and directional
regular expressions with up to five intervening words. Its complete executable
patterns have not been inspected here. The example therefore proves a limit
of the CIR component, not an exploit of the implementation or an equal score
under the whole composite. A banking evaluation should score the proposed
cause, outcome, direction, conditioning information and intervention
separately, while retaining the free-text explanation for review.
\citep{exec17_benchmark_cherif,exec17_benchmark_exact}

The paper selects the composite weights from fifteen candidates using 24
responses rated by two experts. The reported selected Kendall correlation
of 0.78 describes agreement on the selection set. It does not provide a
held-out validation of the selected metric. The required next evidence is a
fixed metric applied to different mechanisms and independently adjudicated
responses, including plausible keyword-rich errors, appropriate refusals and
alternative correct phrasings. This protects against a score that rewards
similar vocabulary while missing the banking decision's meaning.
\citep{exec17_benchmark_cherif}

Reproduction is also a substantive limit. The methods describe code, prompts,
gold annotations and raw results as released publicly, whereas the Data
Availability Statement on physical page 33 says they are available on
reasonable request. No public executable copy has been established by this
inspection. Both statements must be retained; the paper alone cannot qualify
a runnable integration. Its historical zero API charge under a free tier and
roughly eighteen hours of runtime likewise leave annotation, review, waiting,
infrastructure and current prices unmeasured. These limits support using the
paper as a source of test-design ideas, while leaving direct adoption and full
comparative economics unresolved.~\citep{exec17_benchmark_cherif}

\subsection{The monetary meaning of answering and withholding}

An all-No benchmark rewards withholding regardless of its opportunity cost.
To expose that omission, consider a deliberately stipulated banking review
task. A request is either supported or unsupported by its supplied information.
Endorsing a supported request earns USD 20; endorsing an unsupported request
loses USD 80; withholding earns zero. These amounts represent the value and
error cost of a review decision, not the deposit offer's actual causal effect
or an observed bank return. Let $\rho$ be the supported fraction, $t$ the
probability of endorsement given support and $f$ the probability of endorsement
without support. Then
\begin{align}
 A&=\rho t+(1-\rho)(1-f),
       \label{exec17:benchmark-accuracy}\\
 V&=20\rho t-80(1-\rho)f.
       \label{exec17:benchmark-value}
\end{align}
Here $A$ is classification accuracy and $V$ is expected USD per request.
They answer different questions. At the balanced mix $\rho=1/2$, the value
reduces to $10t-40f$.
\citep{exec17_benchmark_exact}

\begin{table}[htbp]
\centering
\caption{Exact illustrative support-review decisions at a balanced case mix.
All rates and payoffs are stipulated.}
\label{exec17:benchmark-bank-table}
\begin{tabular}{lrrrr}
\toprule
Rule & $t$ & $f$ & Accuracy & USD/request\\
\midrule
A & 0.80 & 0.10 & 0.850 & 4.00\\
B & 0.95 & 0.20 & 0.875 & 1.50\\
Always withhold & 0 & 0 & 0.500 & 0\\
\bottomrule
\end{tabular}
\end{table}

Rule B has higher accuracy and lower monetary value than rule A. A separate
400-request integer enumeration, with 200 supported and 200 unsupported
requests, gives A 160 true and 20 false endorsements, and B 190 true and 40
false endorsements. The resulting totals are USD 1,600 and USD 600,
respectively, reproducing the per-request values in the table. Always
withholding earns zero and obtains 50\% balanced accuracy, although it would
obtain 100\% accuracy if every label were unsupported. The comparison explains
why the refusal benchmark's label mix and the banking loss function must
remain visible.
\citep{exec17_benchmark_exact}

For a calibrated probability $p$ that a request is supported, endorsement has
conditional value $20p-80(1-p)$. It beats zero-value withholding only for
$p>0.8$, with indifference at 0.8. Calibration here concerns the support-review
label in its specified population. It does not calibrate an intervention
effect, a confidence interval or an assistant's self-reported confidence.
Changing the case mix or loss changes the decision problem. A public prototype
can vary these declared quantities to expose an unfavorable tradeoff; it
cannot discover the bank's actual losses merely by computing this threshold.
\citep{exec17_benchmark_exact}

\subsection{Repeated prompts are not independent mechanisms}

Prompting strategies should ordinarily be compared on the same mechanisms,
because differences in case difficulty otherwise obscure the contrast. Let
$D_{ij}$ be the paired utility difference between two strategies on mechanism
$i$ and within-mechanism variant $j$. An illustrative random-effects law is
\begin{equation}
 D_{ij}=\mu+U_i+E_{ij},\qquad
 \operatorname{Var}(U_i)=\tau^2,\quad
 \operatorname{Var}(E_{ij})=\sigma^2,
 \label{exec17:benchmark-cluster-law}
\end{equation}
with independent mechanisms, independent idiosyncratic terms and all such
terms mutually independent. Shared $U_i$ makes variants of the same mechanism
dependent even when their generation calls are separate. With $n$ mechanisms
and $r$ equally weighted variants per mechanism, the grand mean has variance
\begin{equation}
 \operatorname{Var}(\overline D)
 =\frac{\tau^2+\sigma^2/r}{n},\qquad
 \frac{\operatorname{Var}(\overline D)}{(\tau^2+\sigma^2)/(nr)}
 =1+(r-1)\frac{\tau^2}{\tau^2+\sigma^2}.
 \label{exec17:benchmark-cluster-variance}
\end{equation}
The denominator in the ratio is the variance that would follow from falsely
treating the variants as independent. For $n=12$, $r=8$ and
$\tau^2=\sigma^2=1$, the true variance is $3/32$, the independent-case
calculation gives $1/48$, and their ratio is $4.5$. Counting every ordered
within-mechanism covariance gives the same result. These are exact deductions
under the stated illustrative law, not estimates of dependence in either
paper. In particular, this calculation does not replace or reconstruct the
published significance tests or bootstrap intervals.
\citep{exec17_benchmark_exact}

Extra variants reduce the idiosyncratic contribution but do not remove the
mechanism contribution. A proposed prototype should therefore preserve the
mechanism identifier, compare strategies in pairs and allocate additional
budget to new mechanisms when mechanism variation dominates. With very few
mechanisms, simply choosing a cluster bootstrap does not guarantee reliable
uncertainty. The design must state its sampling population, weighting and
inferential assumptions, and may need to report a finite-case comparison
without a population claim. None of the twelve source scenarios becomes an
independent bank population because many prompts were generated from it.

These two benchmarks suggest complementary roles in the proposal. CELLO
offers a structured translation from entities through graphs to questions,
together with a warning that refusal polarity can conceal a semantic failure.
CausalReasonBench offers explicit automatic-score definitions and prompting
contrasts, together with limits in token matching, metric selection, repeated
mechanisms and reproducibility. The proposed evaluation should adopt those
test ideas only with explicit labels and information regimes, constant-answer
comparators, semantic adjudication, paired mechanisms and declared review
losses. Public cases and synthetic mechanisms can establish whether a narrow
prototype is worth further investigation. Authentic assignment information,
bank outcomes and institutional economics remain different evidence.
\citep{exec17_benchmark_cello,exec17_benchmark_cherif}
\label{exec17:benchmark-development-end}
